\documentclass[10pt,a4paper]{amsart}
\usepackage[margin=19mm]{geometry}
\usepackage{amsmath,amssymb,mathtools}
\usepackage[scr=rsfs,cal=euler]{mathalfa}
\usepackage{xfrac}
\usepackage[unicode,hidelinks,bookmarks=false]{hyperref}
\newcommand{\ie}{i.e.\ }
\newcommand{\cf}{cf.\ }
\newcommand{\resp}{respectively\ }
\newcommand{\Subsection}{\subsection}
\providecommand{\nobreakdash}{\nobreak\hbox{-}}
\setlength{\emergencystretch}{2em}
\multlinegap0pt
\usepackage{tikz}
\usepackage{amssymb}
\usepackage{mathtools}
\usepackage[shortlabels]{enumitem}
\usepackage{xcolor}
\newtheorem{corollary}{Corollary}
\title{Classification of Extended\\ Abelian Chern--Simons Theories}
\author{Daniel Galviz}
\date{Source: arXiv:2604.02929v1 (CC BY 4.0)}
\begin{document}
\maketitle

\noindent\emph{Source and licence.} Classification of Extended\\ Abelian Chern--Simons Theories. Version arXiv:2604.02929v1 (CC BY 4.0), distributed by arXiv under the Creative Commons Attribution 4.0 International licence, http://creativecommons.org/licenses/by/4.0/, as recorded in the arXiv licence metadata for this exact version and shown on the abstract page https://arxiv.org/abs/2604.02929v1. Full licence notice: \texttt{licenses/arxiv-2604.02929-CC-BY-4.0.txt}.\par\smallskip

\noindent\textit{Editorial scope.} Counted unit: the article's corollary corollary (source0.tex lines 207--214). The article's complete original proof of it is delivered with it (source0.tex lines 216--228, one \texttt{proof} environment of 13 lines). No mathematics is written, repaired or re-derived: the statement and the proof are the article's own lines, in the article's own notation, and no retained line is substituted or re-flowed.  No further source-local context is needed: every cross-reference of every retained line resolves inside the two delivered spans.  Nothing was omitted from any retained span, so the package carries no omission record.  No retained line contains a citation, so the wrapper carries no bibliography and no bibliography entry.  No retained line contains a figure environment, an image-inclusion command or drawing code, so no image is delivered and the package declares no asset dependency.  Editorial formatting only, in addition: an AMS \texttt{amsart} preamble that restates the article's own declarations and macros used by the retained lines, namely the environments \texttt{corollary} on separate counters, together with the standard packages those lines need.  The retained environments print in this document as Corollary 1 in order of appearance, whereas the article prints the same items with the numbering of its own full document.  The complete native source of the article as distributed in the arXiv e-print for this exact version is bundled unchanged as \texttt{sources/197769/source0.tex}.
\begin{corollary}
Let $(\Lambda_1,K_1)$ and $(\Lambda_2,K_2)$ be even, integral, nondegenerate
lattices. If $(G_{K_1},q_{K_1})\cong (G_{K_2},q_{K_2}),$ then
\[
Z^{CS}_{\mathbb T_1,K_1}\cong Z^{CS}_{\mathbb T_2,K_2}
\]
as symmetric monoidal extended $(2+1)$-dimensional TQFTs.
\end{corollary}

\begin{proof}
If $(G_{K_1},q_{K_1})\cong (G_{K_2},q_{K_2})$, then the pointed modular tensor
categories $\mathcal C(G_{K_1},q_{K_1})$ and $
\mathcal C(G_{K_2},q_{K_2})$ are braided ribbon equivalent. Hence their  Reshetikhin--Turaev theories are naturally isomorphic as symmetric monoidal
functors. Applying the extended equivalence theorem to both lattice
presentations yields
\[
Z^{CS}_{\mathbb T_1,K_1}\cong Z^{RT}_{\mathcal C(G_{K_1},q_{K_1})}
\cong Z^{RT}_{\mathcal C(G_{K_2},q_{K_2})}
\cong Z^{CS}_{\mathbb T_2,K_2}\,\,,
\]
which proves the claim.
\end{proof}

\end{document}
